\begin{figure}[tbp]
\centering
\begin{minipage}[t]{.47\linewidth}
\centering\ViewHeading{(a) Nilpotent mark: strict growth}
\begin{tikzpicture}
\node[vseed,label=below:$0$] (zero) at (0,0) {};
\node[vseed,label=below:{$[e_1]$}] (first) at (2.8,0) {};
\node[vpoint,label=above:{$[e_2]$}] (second) at (0,1.6) {};
\node[vpoint,label=above:{$[e_1+e_2]$}] (sum) at (2.8,1.6) {};
\draw[varrow] (first)--(second) node[midway,above right,font=\footnotesize] {$\overline G$};
\draw[varrow] (second)--(zero);
\draw[varrow] (sum)--(second);
\node[align=center] at (1.4,-1.35)
 {$N^{(0)}=\{0,[e_1]\}$\\$N^{(1)}=H=\F_2^2$};
\end{tikzpicture}
\end{minipage}\hfill
\begin{minipage}[t]{.47\linewidth}
\centering\ViewHeading{(b) Zero mark: no growth}
\begin{tikzpicture}
\node[vseed,label=below:$0$] (zero) at (0,0) {};
\node[vseed,label=below:{$[e_1]$}] (first) at (2.8,0) {};
\node[circle,draw=viewgray,inner sep=2.5pt,label=above:{$[e_2]$}] (second) at (0,1.6) {};
\node[circle,draw=viewgray,inner sep=2.5pt,label=above:{$[e_1+e_2]$}] (sum) at (2.8,1.6) {};
\draw[varrow] (first)--(zero);
\draw[varrow] (second)--(zero);
\draw[varrow] (sum)--(zero);
\node[align=center] at (1.4,-1.35)
 {$N^{(m)}=\{0,[e_1]\}$ for all $m$\\$\dim\Nex=1$};
\end{tikzpicture}
\end{minipage}
\ViewHeading{(c) A swap can descend to the identity}
\begin{tikzpicture}
\node[vbox,text width=4.8cm] (swap) at (-3.8,0)
 {$\mathsf Z=\F_2^2$, $G(x,y)=(y,x)$\\$(1,0)\longleftrightarrow(0,1)$};
\node[vresult,text width=5.3cm] (identity) at (3.8,0)
 {$\mathsf B=\langle(1,1)\rangle$\\$[(1,0)]=[(0,1)]$,\quad $\overline G=\id$};
\draw[varrow] (swap)--(identity) node[midway,above,font=\footnotesize] {quotient};
\end{tikzpicture}
\caption{Exact finite-field witnesses, drawn as finite points rather
than continuous planes. Panels (a) and (b) use the same quotient of
$\F_2^3$ by $\langle e_3\rangle$ and the same seed; only the mark changes.
Orange points are the initial seed subspace, blue points are added by
closure, and hollow points are outside the response. Zero is fixed in
both panels (its self-loop is omitted). In (a), linear span adds
$[e_1+e_2]$ as well as the new direction. Panel (c) already has $V=H$.}
\label{fig:response-witnesses}
\end{figure}